\section{Biện pháp thứ nhất: chặn độ dài bước}
\label{sec:ch03-clipping}

\index{clipping!thuật toán}%
Toàn bộ biện pháp nằm gọn trong bốn dòng giả mã dưới đây, chép theo thuật toán 1
của bài.

\begin{algorithm}[htbp]
  \caption{Clipping theo norm, giả mã theo thuật toán 1 của bài báo}
  \label{alg:ch03-clipping}
  \begin{algorithmic}[1]
    \State \(\hat{\vect{g}} \gets \pd{L}{\theta}\)
    \If{\(\norm{\hat{\vect{g}}} \ge \tau\)}
      \State \(\hat{\vect{g}} \gets \dfrac{\tau}{\norm{\hat{\vect{g}}}}\, \hat{\vect{g}}\)
    \EndIf
  \end{algorithmic}
\end{algorithm}

\(\tau\) là ngưỡng. Phép gán ở dòng ba là một phép co tỷ lệ, nên hướng của
\(\hat{\vect{g}}\) không đổi và bước đi vẫn là hướng giảm đối với mini-batch
hiện tại. Chỉ độ dài bị chặn lại ở \(\tau\).

Ba tính chất của giả mã này đáng giữ trong đầu. Nó chỉ đọc thứ mà backward pass
đã tính xong, nên không tốn thêm gì. Nó so sánh với một số cố định, nên trên
vùng phẳng nó không làm gì cả. Và nó dài bốn dòng.

Bài ghi rõ xuất xứ: thuật toán này rất gần với đề xuất của Tomas Mikolov, người
cắt theo từng phần tử thay vì theo norm, và bài chỉ lệch khỏi bản gốc để bảo
đảm bước đi luôn là hướng giảm đối với mini-batch hiện tại. Họ nói thêm rằng
trên thực tế hai biến thể hành xử tương tự nhau.

\subsection*{Một bước tại đỉnh vách}

Tôi lấy đúng điểm dốc nhất ở mục~\ref{sec:ch03-he-dong-luc}, tìm lại bằng chính
hàm mà mục đó gọi, learning rate 0.1, rồi đi một bước hai lần: một lần không
cắt, một lần có cắt.

\begin{measured}{một bước gradient descent tại \(b = -2.6123413600\)}
\begin{minted}[fontsize=\footnotesize]{text}
cost at start      0.14960413
dE/dw              -3.577721e+03
dE/db              -5.751099e+03
||grad||           6.773125e+03

 thresh  fired        step        w after        b after        cost
   none  False  6.7731e+02   362.77214318   572.49753848  0.09000000
    1.0   True  1.0000e-01     5.05282231    -2.52743080  0.02586525
    0.1   True  1.0000e-02     5.00528223    -2.60385030  0.01215804
\end{minted}

Đo bằng \code{experiments/ch03\_clipping.py} tại tag \code{ch03}.
\end{measured}

\index{clipping!hiệu quả đo được}%
Bước không cắt đưa \((w, b)\) từ \((5.0, -2.61)\) tới
\((362.77214318,\ 572.49753848)\). Không phải một bước hơi dài. Nó rời khỏi
cửa sổ mà \tn{mặt lỗi}{error surface} được vẽ, rời khỏi mọi vùng tham số có nghĩa, và đi xa hơn
điểm xuất phát khoảng sáu trăm bảy mươi bảy đơn vị trong một không gian mà toàn
bộ bài toán nằm gọn trong một hình vuông cạnh chưa tới một.

Chi phí sau bước đó là 0.09000000, và con số tròn trịa ấy không phải trùng hợp.
Với \(w\) và \(b\) lớn như vậy, \(\sigma(a_{50})\) bão hòa về đúng 1 trong
float64, nên chi phí bằng \((1 - 0.7)^2\). Mạng không rơi vào một cực tiểu tệ
hơn. Nó rời khỏi bài toán, tới một chỗ mà gradient bằng không và không có đường
quay lại.

Hai bước có cắt thì ở lại. Ngưỡng 1.0 đưa tới \((5.05282231,\ -2.52743080)\) và
chi phí xuống 0.02586525. Ngưỡng 0.1 đưa tới \((5.00528223,\ -2.60385030)\) và
chi phí xuống 0.01215804. Cả hai đều giảm chi phí so với 0.14960413 lúc xuất
phát, còn bước không cắt thì tăng.

Tỷ số độ dài giữa bước không cắt và bước cắt ở ngưỡng 1.0 là 6.7731 nhân mười
mũ ba, tức 3.8308 bậc. Đó là toàn bộ khoảng cách giữa một lần huấn luyện hỏng
và một lần huấn luyện chạy tiếp.

\subsection*{Chọn ngưỡng, và tìm nó ở đâu trong PyTorch}

\index{clipping!chọn ngưỡng}%
Bài đưa một heuristic thay vì một công thức: nhìn thống kê chuẩn gradient trung
bình qua đủ nhiều bước cập nhật rồi lấy ngưỡng quanh đó. Họ nói thêm rằng huấn
luyện không nhạy lắm với siêu tham số này, và thuật toán chạy tốt cả khi ngưỡng
khá nhỏ. Số đo phía trên ủng hộ điều đó: ngưỡng 1.0 và ngưỡng 0.1 lệch nhau mười
lần mà cả hai đều làm đúng việc cần làm, chỉ khác nhau ở chỗ đi được bao xa.

Bạn không phải tự viết thuật toán~\ref{alg:ch03-clipping}. PyTorch có sẵn, và
repo companion gói lại đúng một câu lệnh ấy:

\begin{minted}{python}
def clip_parameters(parameters, threshold: float) -> float:
    return torch.nn.utils.clip_grad_norm_(
        parameters, max_norm=threshold
    ).item()
\end{minted}

Hàm của thư viện co tỷ lệ trên toàn bộ danh sách tham số coi như một vector duy
nhất, sửa \code{.grad} tại chỗ, và trả về chuẩn nó thấy trước khi cắt. Gọi nó
sau \code{loss.backward()} và trước \code{optimizer.step()}.

Một chi tiết tôi chỉ thấy khi viết test cho nó: hàm này \emph{không} phải thuật
toán~\ref{alg:ch03-clipping}. Nó nhân với \(\tau / (\norm{\hat{\vect{g}}} + 10^{-6})\)
chứ không phải \(\tau / \norm{\hat{\vect{g}}}\), để khỏi chia cho không. Cái giá
là kết quả nằm hơi dưới ngưỡng: với một gradient có chuẩn \(\sqrt{50}\), chuẩn
sau khi cắt là 0.9999998585786636 chứ không phải 1. Không đáng kể trong huấn
luyện, nhưng đáng biết nếu bạn định khẳng định một bất biến về chuẩn sau khi
cắt.
